\documentclass{exam}
\usepackage[utf8]{inputenc}
\usepackage{color}
\usepackage{amsmath}
\usepackage[italian]{babel}
\usepackage{comment}
\usepackage{graphicx}
\usepackage[siunitx]{circuitikz}
\newcommand{\splitcell}[2][c]{%
  \begin{tabular}[c]{@{}c@{}}\strut#2\strut\end{tabular}%
}
\usepackage{natbib}
\usepackage{listings}
\usepackage{xparse}
\usepackage{hyperref}
\usepackage{float}
\usepackage{tcolorbox}

\NewDocumentCommand{\codeword}{v}{%
\texttt{\textcolor{blue}{#1}}%
}
\usepackage{pgfplots}
\pgfplotsset{compat = newest}
\usepackage{chngcntr}
\counterwithin{table}{section}
\counterwithin{figure}{section}

\usepackage{listings}
\usepackage{xcolor}

\definecolor{codegreen}{rgb}{0,0.6,0}
\definecolor{codegray}{rgb}{0.5,0.5,0.5}
\definecolor{codepurple}{rgb}{0.58,0,0.82}
\definecolor{backcolour}{rgb}{0.95,0.95,0.92}

\lstdefinestyle{mystyle}{
    backgroundcolor=\color{backcolour},   
    commentstyle=\color{codegreen},
    keywordstyle=\color{magenta},
    numberstyle=\tiny\color{codegray},
    stringstyle=\color{codepurple},
    basicstyle=\ttfamily\footnotesize,
    breakatwhitespace=false,         
    breaklines=true,                 
    captionpos=b,                    
    keepspaces=true,                 
    numbers=left,                    
    numbersep=5pt,                  
    showspaces=false,                
    showstringspaces=false,
    showtabs=false,                  
    tabsize=2
}

\lstset{style=mystyle}

\begin{document}

\newcommand{\Exjobbsnummer}[1]{
	\begin{tikzpicture}[overlay, remember picture]
		\path (current page.north east) ++(-1,-1) node[below left] {{\small #1}};
	\end{tikzpicture}
}

\newcommand{\Examensjobbspoang}[1]{
	\begin{tikzpicture}[overlay, remember picture]
		\path (current page.north east) ++(-1,-1.5) node[below left] {{\normalsize \scshape Examensarbete #1 HP}};
	\end{tikzpicture}
}

\newcommand{\datum}[1]{
	\begin{tikzpicture}[overlay, remember picture]
		\path (current page.north east) ++(-1,-2.0) node[below left] {{\normalsize #1}};
	\end{tikzpicture}}

\newcommand{\storlitentitel}[2]{
\center
\rule[0.2cm]{13cm}{0.1cm}
{ \huge \bfseries #1}\\[0.4cm] % Title of your document
{\Large \slshape #2}\\[0.4cm]
\rule[0.2cm]{13cm}{0.1cm}\\[3cm]

}

\newcommand{\Namn}[2]{
	\begin{minipage}{0.5\textwidth}
		\normalsize
		\centering
		#1 \textsc{#2}\\
	\end{minipage}\\
}

\newcommand{\LoggaSwe}{
	\textsc{\Huge Laboratorio \\[0.3cm] di Internet}\\[0.7cm]
	\includegraphics[scale=.06]{polito_logo_2021_blu.jpg}\\[1.5cm]
}

\newcommand{\LoggaEng}{
	\textsc{\Huge Uppsala University}\\[0.7cm]
	\includegraphics[scale=.1]{Uppsala_University_seal_svg.png}\\[0.5cm]
}


% -----------------------------------------------
%           HÄR BÖRJAR TITELSIDAN
%------------------------------------------------
\begin{titlepage}

	\center

	%-------------------------------------------------
	%	INFORMATION ATT FYLLA I
	%-------------------------------------------------
	\Exjobbsnummer{Anno Accademico 2021/2022}
	%\datum{2021/2022}

	\LoggaSwe
	% \LoggaEng - Byt till engelska


	\storlitentitel{\\Report II}{Stima di velocità tramite ping}

	\Large Gruppo 14\\
	\Namn{guest1}{(s000000)}
	\Namn{Alessandro Ciullo}{(s310023)}
	\Namn{guest3}{(s000000)}



	\vfill

\end{titlepage}
\pagebreak


\section{Network Configuration}
\begin{table}
  \caption{Net hosts}\label{tab:net-hosts}
  \begin{center}
    \begin{tabular}{|l|l|}
      \hline
      \textbf{Host name} & \textbf{IP address} \\
      \hline
      H1 & 10.0.3.1 \\
      \hline
      H2 & 10.0.3.2 \\
      \hline
      H3 & 10.1.3.1 \\
      \hline
      H4 & 10.1.3.2 \\
      \hline
    \end{tabular}
  \end{center}
\end{table}

Test1
H1 mtu 1500, H2 mtu 1500
\verb|H2$ ping H1 -s 2000 -c 1 -M want|

Test2
H1 mtu 1500, H2 mtu 1500
\verb|H2$ ping H1 -s 2000 -c 1 -M do|
H2 blocks the packet at IP level, the packet is not event sent through the network card

Test3
H1 mtu 1500, H2 mtu 1500
\verb|H2$ ping H1 -s 2000 -c 1 -M dont|
Same result as test1

Test4
H1 mtu 800, H2 mtu 1500
\verb|H2$ ping H1 -s 2000 -c 1 -M want|

Test6
H1 mtu 800, H2 mtu 1500
\verb|H2$ ping H1 -s 2000 -c 1 -M dont|

Test7
H1 mtu 800, H2 mtu 1500, H3 mtu 1200, H4 mtu 900
\verb|H2$ ping H4 -s 1300 -c 2 -M want|

Test8
H1 mtu 800, H2 mtu 1500, H3 mtu 1200, H4 mtu 900
\verb|H2$ ping H4 -s 1300 -c 2 -M do|

Test9
H1 mtu 800, H2 mtu 1500, H3 mtu 1200, H4 mtu 900
\verb|H2$ ping H4 -s 1300 -c 2 -M dont|

Test10
H1 mtu 1500, H2 mtu 9000
\verb|H2$ ping H1 -s 65000 -c 1|

Test11
H1 mtu 1501, H2 mtu 9000
\verb|H2$ ping H4 -s 4182 -c 1|

Test12
H1 mtu 1501, H2 mtu 9000
\verb|H2$ ping H4 -s 4183 -c 1|
From H2 wireshark -> ip is sent in the network (no response found!)



\end{document}
